% !PRINTED 327
\weberpage{327}

%<<CONT
If $x$ lies near a root, then one of the two values
of $X$ gives a still more accurate value of this root.

\begin{center}
\rule{2cm}{0.4pt}
\end{center}

Just as an interval is determined by $X_1, X_2$ in which
no root of the equation lies at all, so one can also determine an
interval in which roots certainly lie. Here too
one best starts from the homogeneous form and makes use of
the linear transformation.

Let $f(x, y)$ then be, as before, a form of degree $n$ with none but
real roots; and let further $x:y$, $\xi:\eta$, $\xi':\eta'$ be three real
values, of which two are arbitrary, and of which the
third, through the relation

\weq{16}{\begin{split}
{}[\xi f'(x) + \eta f'(y)]\,[\xi' f'(x) + \eta' f'(y)] \\
{}- H(\xi y - \eta x)\,(\xi' y - \eta' x) = 0
\end{split}}

depends on the two others. The equation (16) remains
unchanged if we subject the three pairs of quantities $x, y$; $\xi, \eta$; $\xi', \eta'$
simultaneously to the same linear transformation. Hence
all conclusions that we draw from a special form of this
equation can be generalised. The linear substitution
can be so determined that for the new variables $y = 0$,
$\xi = 0$; then (16) yields for $\xi'$, $\eta'$ the equation

\[
a_1(n\,a_0\,\xi' + a_1\,\eta') = [2\,n\,a_0\,a_2 - (n-1)\,a_1^2]\eta',
\]

or, if one assumes $a_0 = 1$, $\eta' = 1$,

\[
\xi' = -\frac{a_1^2 - 2\,a_2}{a_1}.
\]

But if, as above, we denote by $\alpha$ the roots of $f$,
we have

\[
a_1^2 - 2\,a_2 = S(\alpha^2), \quad a_1 = -S(\alpha);
\]

and thus

\[
\xi' = \frac{S(\alpha^2)}{S(\alpha)}.
\]

For this we may also write:

\[
S[\alpha(\alpha - \xi')] = 0,
\]

and from this it follows, since all the $\alpha$ are to be real, that the
products $\alpha(\alpha - \xi')$ can be neither all positive nor all negative,
or in other words, one part of the roots must
%CONT>>
